% GENERATED FILE. Edit canonical lessons, then run npm run generate:books.
% canonical-source-hash: fnv1a64:d008a523149bd374
% canonical-lessons: TE-C517-laksanam, TE-C517-maikam, TE-C517-ayasapadu, TE-C517-kakku, TE-C517-killanoppulu, TE-R517-review-510-517

\chapter{Symptoms}
\label{ch:te-symptoms}
\hlchaptermodality{517}

\begin{chapteropening}
\textbf{By the end of this chapter:} \emph{I can describe symptoms such as dizziness, breathlessness and joint pain.}

Complete the last lesson of chapter 517: I can describe symptoms such as dizziness, breathlessness and joint pain.
\end{chapteropening}

\section[\texorpdfstring{\te{లక్షణం} (lakṣaṇaṁ)}{లక్షణం (lakṣaṇaṁ)}]{\te{లక్షణం} (lakṣaṇaṁ) — a symptom; a trait}

\label{lesson:TE-C517-laksanam}



\begin{quote}
Before the new one: say the Telugu for an allergy, then the Telugu for sleeplessness.
\end{quote}

\subsection*{You'll want to know: \te{లక్షణం}}
\textbf{\te{లక్షణం}} — \emph{lakṣaṇaṁ} — \textquotedblleft{}a symptom; a trait\textquotedblright{}.

\begin{quote}
\te{జ్వరం}, \te{దగ్గు} \te{ఆ} \te{జబ్బు} \te{లక్షణాలు}.
\end{quote}

\begin{quote}
\emph{jvaraṁ, daggu ā jabbu lakṣaṇālu.}
\end{quote}

\begin{quote}
Fever and cough are the symptoms of that illness.
\end{quote}

Plural: \te{లక్షణాలు}, \emph{lakṣaṇālu}.

Also a person's traits: \te{మంచి} \te{లక్షణాలు}, good qualities.

\begin{cousinweb}
\textbf{\te{లక్షణం}} (\emph{lakṣaṇaṁ}) is borrowed from Sanskrit \emph{lakṣaṇa}, \textquotedblleft{}line, mark\textquotedblright{} (CDIAL 10882): the marks by which an illness or a person is known.
\end{cousinweb}

\begin{grammarlens}[title={What you've built}]
One more word for describing what you lived and felt.
\end{grammarlens}

\subsection*{Your turn}
\begin{itemize}
\raggedright
  \item \emph{Say it:} \emph{lakṣaṇaṁ}
  \item \emph{Say it:} \emph{lakṣaṇaṁ}, once more
  \item \emph{Say it:} \emph{jvaraṁ, daggu ā jabbu lakṣaṇālu.} — Fever and cough are the symptoms of that illness.
  \item \emph{Recall:} say the Telugu for an allergy, then the Telugu for sleeplessness, then say \emph{lakṣaṇaṁ} again
\end{itemize}

\begin{tcolorbox}[breakable,colback=teal!4,colframe=teal!35!black,title={Before you move on}]
What does \te{లక్షణం} mean? (\textquotedblleft{}A symptom; a trait\textquotedblright{}.) Say it once more.
\end{tcolorbox}
\section[\texorpdfstring{\te{మైకం} (maikaṁ)}{మైకం (maikaṁ)}]{\te{మైకం} (maikaṁ) — dizziness; a daze}

\label{lesson:TE-C517-maikam}



\begin{quote}
Before the new one: say the Telugu for sleeplessness, then the Telugu for a symptom.
\end{quote}

\subsection*{You'll want to know: \te{మైకం}}
\textbf{\te{మైకం}} — \emph{maikaṁ} — \textquotedblleft{}dizziness; a daze\textquotedblright{}.

\begin{quote}
\te{ఎండలో} \te{నడిచాక} \te{నాకు} \te{మైకంగా} \te{అనిపించింది}.
\end{quote}

\begin{quote}
\emph{eṇḍalō naḍicāka nāku maikaṅgā anipiñcindi.}
\end{quote}

\begin{quote}
After walking in the sun, I felt giddy.
\end{quote}

\te{మైకంగా} \te{ఉంది}, I feel giddy; also the daze of sleep.

\begin{cousinweb}
\textbf{\te{మైకం}} (\emph{maikaṁ}) is native Dravidian: DEDR 4706 lists Telugu \emph{maikamu}, \textquotedblleft{}intoxication, unconsciousness, drowsiness\textquotedblright{}, in the article of Tamil \emph{mayakku}, \textquotedblleft{}to bewilder, to confuse\textquotedblright{}.
\end{cousinweb}

\begin{grammarlens}[title={What you've built}]
One more word for describing what you lived and felt.
\end{grammarlens}

\subsection*{Your turn}
\begin{itemize}
\raggedright
  \item \emph{Say it:} \emph{maikaṁ}
  \item \emph{Say it:} \emph{maikaṁ}, once more
  \item \emph{Say it:} \emph{eṇḍalō naḍicāka nāku maikaṅgā anipiñcindi.} — After walking in the sun, I felt giddy.
  \item \emph{Recall:} say the Telugu for sleeplessness, then the Telugu for a symptom, then say \emph{maikaṁ} again
\end{itemize}

\begin{tcolorbox}[breakable,colback=teal!4,colframe=teal!35!black,title={Before you move on}]
What does \te{మైకం} mean? (\textquotedblleft{}Dizziness; a daze\textquotedblright{}.) Say it once more.
\end{tcolorbox}
\section[\texorpdfstring{\te{ఆయాసపడు} (āyāsapaḍu)}{ఆయాసపడు (āyāsapaḍu)}]{\te{ఆయాసపడు} (āyāsapaḍu) — to pant; to get breathless}

\label{lesson:TE-C517-ayasapadu}



\begin{quote}
Before the new one: say the Telugu for a symptom, then the Telugu for dizziness.
\end{quote}

\subsection*{You'll want to know: \te{ఆయాసపడు}}
\textbf{\te{ఆయాసపడు}} — \emph{āyāsapaḍu} — \textquotedblleft{}to pant; to get breathless\textquotedblright{}.

\begin{quote}
\te{కొండ} \te{ఎక్కి} \te{తాతయ్య} \te{ఆయాసపడ్డారు}.
\end{quote}

\begin{quote}
\emph{koṇḍa ekki tātayya āyāsapaḍḍāru.}
\end{quote}

\begin{quote}
Climbing the hill, grandpa got breathless.
\end{quote}

Past (I …): \te{ఆయాసపడ్డాను}, \emph{āyāsapaḍḍānu}.

\te{ఆయాసం}, breathlessness, + \te{పడు}.

\begin{cousinweb}
\textbf{\te{ఆయాసపడు}} (\emph{āyāsapaḍu}) is built from \te{ఆయాసం}, from Sanskrit \emph{āyāsa}, \textquotedblleft{}effort, exertion; fatigue\textquotedblright{} (Monier-Williams) and the native verb \emph{paḍu}, \textquotedblleft{}to feel, to suffer\textquotedblright{} (DEDR 3853).
\end{cousinweb}

\begin{grammarlens}[title={What you've built}]
One more word for describing what you lived and felt.
\end{grammarlens}

\subsection*{Your turn}
\begin{itemize}
\raggedright
  \item \emph{Say it:} \emph{āyāsapaḍu}
  \item \emph{Say it:} \emph{āyāsapaḍu}, once more
  \item \emph{Say it:} \emph{koṇḍa ekki tātayya āyāsapaḍḍāru.} — Climbing the hill, grandpa got breathless.
  \item \emph{Recall:} say the Telugu for a symptom, then the Telugu for dizziness, then say \emph{āyāsapaḍu} again
\end{itemize}

\begin{tcolorbox}[breakable,colback=teal!4,colframe=teal!35!black,title={Before you move on}]
What does \te{ఆయాసపడు} mean? (\textquotedblleft{}To pant; to get breathless\textquotedblright{}.) Say it once more.
\end{tcolorbox}
\section[\texorpdfstring{\te{కక్కు} (kakku)}{కక్కు (kakku)}]{\te{కక్కు} (kakku) — to throw up; to spit out}

\label{lesson:TE-C517-kakku}



\begin{quote}
Before the new one: say the Telugu for dizziness, then the Telugu for to pant.
\end{quote}

\subsection*{You'll want to know: \te{కక్కు}}
\textbf{\te{కక్కు}} — \emph{kakku} — \textquotedblleft{}to throw up; to spit out\textquotedblright{}.

\begin{quote}
\te{బస్సులో} \te{ప్రయాణం} \te{చేస్తూ} \te{పాప} \te{కక్కింది}.
\end{quote}

\begin{quote}
\emph{bassulō prayāṇaṁ cēstū pāpa kakkindi.}
\end{quote}

\begin{quote}
The little girl threw up on the bus journey.
\end{quote}

Past (I …): \te{కక్కాను}, \emph{kakkānu}.

Colloquial; \te{వాంతి} \te{చేసుకో} is the polite way to say it.

\begin{cousinweb}
\textbf{\te{కక్కు}} (\emph{kakku}) is native Dravidian: DEDR 1079 lists Telugu \emph{k(r)akku}, \textquotedblleft{}to vomit, to disgorge\textquotedblright{}.
\end{cousinweb}

\begin{grammarlens}[title={What you've built}]
One more word for describing what you lived and felt.
\end{grammarlens}

\subsection*{Your turn}
\begin{itemize}
\raggedright
  \item \emph{Say it:} \emph{kakku}
  \item \emph{Say it:} \emph{kakku}, once more
  \item \emph{Say it:} \emph{bassulō prayāṇaṁ cēstū pāpa kakkindi.} — The little girl threw up on the bus journey.
  \item \emph{Recall:} say the Telugu for dizziness, then the Telugu for to pant, then say \emph{kakku} again
\end{itemize}

\begin{tcolorbox}[breakable,colback=teal!4,colframe=teal!35!black,title={Before you move on}]
What does \te{కక్కు} mean? (\textquotedblleft{}To throw up; to spit out\textquotedblright{}.) Say it once more.
\end{tcolorbox}
\section[\texorpdfstring{\te{కీళ్ళనొప్పులు} (kīḷḷanoppulu)}{కీళ్ళనొప్పులు (kīḷḷanoppulu)}]{\te{కీళ్ళనొప్పులు} (kīḷḷanoppulu) — joint pains}

\label{lesson:TE-C517-killanoppulu}



\begin{quote}
Before the new one: say the Telugu for to pant, then the Telugu for to throw up.
\end{quote}

\subsection*{You'll want to know: \te{కీళ్ళనొప్పులు}}
\textbf{\te{కీళ్ళనొప్పులు}} — \emph{kīḷḷanoppulu} — \textquotedblleft{}joint pains\textquotedblright{}.

\begin{quote}
\te{వర్షాకాలంలో} \te{అమ్మమ్మకు} \te{కీళ్ళనొప్పులు} \te{ఎక్కువ}.
\end{quote}

\begin{quote}
\emph{varṣākālaṁlō ammammaku kīḷḷanoppulu ekkuva.}
\end{quote}

\begin{quote}
In the rainy season grandma's joints ache more.
\end{quote}

\te{కీలు}, a joint (plural \te{కీళ్ళు}), + \te{నొప్పులు}, pains.

\begin{cousinweb}
\textbf{\te{కీళ్ళనొప్పులు}} (\emph{kīḷḷanoppulu}) is built from \te{కీళ్ళు}, \textquotedblleft{}joints\textquotedblright{}, and native \emph{noppulu}, \textquotedblleft{}pains\textquotedblright{} (DEDR 3793).
\end{cousinweb}

\begin{grammarlens}[title={What you've built}]
One more word for describing what you lived and felt.
\end{grammarlens}

\subsection*{Your turn}
\begin{itemize}
\raggedright
  \item \emph{Say it:} \emph{kīḷḷanoppulu}
  \item \emph{Say it:} \emph{kīḷḷanoppulu}, once more
  \item \emph{Say it:} \emph{varṣākālaṁlō ammammaku kīḷḷanoppulu ekkuva.} — In the rainy season grandma's joints ache more.
  \item \emph{Recall:} say the Telugu for to pant, then the Telugu for to throw up, then say \emph{kīḷḷanoppulu} again
\end{itemize}

\begin{tcolorbox}[breakable,colback=teal!4,colframe=teal!35!black,title={Before you move on}]
What does \te{కీళ్ళనొప్పులు} mean? (\textquotedblleft{}Joint pains\textquotedblright{}.) Say it once more.
\end{tcolorbox}
\section[(dialogue)]{Words from chapters 510–517 — a review}

\label{lesson:TE-R517-review-510-517}



\begin{quote}
Say the Telugu for selfishness and negligence.
\end{quote}

\subsection*{Your turn}
Say these aloud:

\begin{itemize}
\raggedright
  \item selfishness, negligence, discipline, self-confidence, curiosity
  \item wits, to be confused, to think back on, to go over again, to remind
  \item a sickness, to fall ill, a slight indisposition, to be bedridden, to come round
  \item a disorder, a contagious disease, a heart attack, paralysis, asthma
  \item an earache, a toothache, a backache, to groan, to faint
  \item the brain, a nerve, the liver, the spine, breath
  \item eyesight, hearing, deafness, an allergy, sleeplessness
  \item a symptom, dizziness, to pant, to throw up, joint pains
\end{itemize}

\begin{tcolorbox}[breakable,colback=teal!4,colframe=teal!35!black,title={Before you move on}]
Selfishness? (\textbf{\te{స్వార్థం}}, \emph{svārthaṁ}.) Negligence? (\textbf{\te{నిర్లక్ష్యం}}, \emph{nirlakṣyaṁ}.) Discipline? (\textbf{\te{క్రమశిక్షణ}}, \emph{kramaśikṣaṇa}.) Self-confidence? (\textbf{\te{ఆత్మవిశ్వాసం}}, \emph{ātmaviśvāsaṁ}.) Curiosity? (\textbf{\te{కుతూహలం}}, \emph{kutūhalaṁ}.) Wits? (\textbf{\te{తెలివితేటలు}}, \emph{telivitēṭalu}.) To be confused? (\textbf{\te{తికమకపడు}}, \emph{tikamakapaḍu}.) To think back on? (\textbf{\te{తలచుకో}}, \emph{talacukō}.) To go over again? (\textbf{\te{నెమరువేయు}}, \emph{nemaruvēyu}.) To remind? (\textbf{\te{గుర్తుచేయు}}, \emph{gurtucēyu}.) A sickness? (\textbf{\te{జబ్బు}}, \emph{jabbu}.) To fall ill? (\textbf{\te{జబ్బుపడు}}, \emph{jabbupaḍu}.) A slight indisposition? (\textbf{\te{నలత}}, \emph{nalata}.) To be bedridden? (\textbf{\te{మంచం} \te{పట్టు}}, \emph{mañcaṁ paṭṭu}.) To come round? (\textbf{\te{తేరుకో}}, \emph{tērukō}.) A disorder? (\textbf{\te{వ్యాధి}}, \emph{vyādhi}.) A contagious disease? (\textbf{\te{అంటువ్యాధి}}, \emph{aṇṭuvyādhi}.) A heart attack? (\textbf{\te{గుండెపోటు}}, \emph{guṇḍepōṭu}.) Paralysis? (\textbf{\te{పక్షవాతం}}, \emph{pakṣavātaṁ}.) Asthma? (\textbf{\te{ఉబ్బసం}}, \emph{ubbasaṁ}.) An earache? (\textbf{\te{చెవిపోటు}}, \emph{cevipōṭu}.) A toothache? (\textbf{\te{పంటినొప్పి}}, \emph{paṇṭinoppi}.) A backache? (\textbf{\te{వెన్నునొప్పి}}, \emph{vennunoppi}.) To groan? (\textbf{\te{మూలుగు}}, \emph{mūlugu}.) To faint? (\textbf{\te{మూర్ఛపోవు}}, \emph{mūrchapōvu}.) The brain? (\textbf{\te{మెదడు}}, \emph{medaḍu}.) A nerve? (\textbf{\te{నరం}}, \emph{naraṁ}.) The liver? (\textbf{\te{కాలేయం}}, \emph{kālēyaṁ}.) The spine? (\textbf{\te{వెన్నెముక}}, \emph{vennemuka}.) Breath? (\textbf{\te{శ్వాస}}, \emph{śvāsa}.) Eyesight? (\textbf{\te{చూపు}}, \emph{cūpu}.) Hearing? (\textbf{\te{వినికిడి}}, \emph{vinikiḍi}.) Deafness? (\textbf{\te{చెవుడు}}, \emph{cevuḍu}.) An allergy? (\textbf{\te{అలెర్జీ}}, \emph{alerjī}.) Sleeplessness? (\textbf{\te{నిద్రలేమి}}, \emph{nidralēmi}.) A symptom? (\textbf{\te{లక్షణం}}, \emph{lakṣaṇaṁ}.) Dizziness? (\textbf{\te{మైకం}}, \emph{maikaṁ}.) To pant? (\textbf{\te{ఆయాసపడు}}, \emph{āyāsapaḍu}.) To throw up? (\textbf{\te{కక్కు}}, \emph{kakku}.) Joint pains? (\textbf{\te{కీళ్ళనొప్పులు}}, \emph{kīḷḷanoppulu}.)
\end{tcolorbox}
